% =============================================================================
% Computing and AI Ethics — OER Conference 2026
% Participant Handout: Key Questions for Cross-Curricular Teaching
% =============================================================================
\documentclass[10pt,letterpaper]{article}

%%% Page geometry
\usepackage[top=0.52in,bottom=0.52in,left=0.48in,right=0.48in,columnsep=0.21in]{geometry}

%%% Core packages
\usepackage{multicol}
\usepackage{tcolorbox}
\tcbuselibrary{skins,breakable}
\usepackage{xcolor}
\usepackage[sfdefault]{plex-sans}
\usepackage[T1]{fontenc}
\usepackage{fontawesome5}
\usepackage{qrcode}
\usepackage{microtype}
\usepackage{enumitem}
\usepackage[hidelinks]{hyperref}
\usepackage{fancyhdr}
\usepackage{booktabs}
\usepackage{array}

\setlength{\parindent}{0pt}
\setlength{\parskip}{0pt}

% ---------------------------------------------------------------------------
% Color palette — matches the course materials
% ---------------------------------------------------------------------------
\definecolor{mPrimary}{HTML}{1E3A5F}   % deep navy
\definecolor{mAccent}{HTML}{D4880A}    % warm amber
\definecolor{tagIT}{HTML}{1A4A7A}
\definecolor{tagSCI}{HTML}{1F6B3A}
\definecolor{tagHLTH}{HTML}{A02020}
\definecolor{tagBUS}{HTML}{1B5E70}
\definecolor{tagLA}{HTML}{5E3085}
\definecolor{tagEDU}{HTML}{8A5900}

% ---------------------------------------------------------------------------
% Inline discipline tag  \dtag{color}{LABEL}
% ---------------------------------------------------------------------------
\newcommand{\dtag}[2]{%
  \raisebox{0.7pt}{%
    \colorbox{#1}{\fontsize{5.4}{6.2}\selectfont
      \textcolor{white}{\textsf{\textbf{\,#2\,}}}}%
  }%
}

% ---------------------------------------------------------------------------
% Chapter card box style
% ---------------------------------------------------------------------------
\tcbset{
  chap/.style={
    enhanced,
    colback=white,
    colframe=mPrimary,
    colbacktitle=mPrimary,
    coltitle=white,
    fonttitle=\footnotesize\sffamily\bfseries,
    boxrule=0.5pt,
    arc=2.5pt,
    left=4pt, right=4pt,
    top=2pt, bottom=2pt,
    toptitle=1.5pt, bottomtitle=1.5pt,
    before skip=2.5pt,
    after skip=2.5pt,
    breakable=false,
  }
}

% ---------------------------------------------------------------------------
% Page footer
% ---------------------------------------------------------------------------
\pagestyle{fancy}
\fancyhf{}
\renewcommand{\headrulewidth}{0pt}
\renewcommand{\footrulewidth}{0.3pt}
\fancyfoot[L]{\sffamily\scriptsize\textcolor{mPrimary}{%
  \textbf{Computing and AI Ethics} \textnormal{---}
  OER Conference 2026 \textnormal{---} Brendan Shea, PhD,
  Rochester Community \& Technical College}}
\fancyfoot[R]{\sffamily\scriptsize\textcolor{mPrimary}{%
  brendanpshea.github.io/computing\_ai\_ethics/ \quad CC BY-NC-SA 4.0}}

% ===========================================================================
\begin{document}
% ===========================================================================

% ---------------------------------------------------------------------------
%  HEADER BANNER
% ---------------------------------------------------------------------------
\begin{tcolorbox}[
  enhanced,
  colback=mPrimary,
  colframe=mPrimary,
  boxrule=0pt,
  arc=3pt,
  left=10pt, right=10pt,
  top=9pt, bottom=9pt,
  before skip=0pt,
  after skip=6pt,
]
\begin{minipage}[c]{0.54\linewidth}
  {\sffamily\bfseries\Large\textcolor{white}{Computing and AI Ethics}}\\[4pt]
  {\sffamily\large\textcolor{mAccent}{%
    Key Questions: Teaching AI Ethics Across the Curriculum}}\\[3pt]
  {\sffamily\small\textcolor{white!70!mPrimary}{%
    Conversation starters and lesson ideas for courses in any discipline}}
\end{minipage}%
\hfill
\begin{minipage}[c]{0.26\linewidth}
  \raggedright
  {\sffamily\footnotesize\textcolor{white!80!mPrimary}{%
    \textbf{Brendan Shea, PhD}\\
    Rochester Community \& Technical College\\
    Philosophy \& Computer Science\\[4pt]
    \faGlobe\enspace\texttt{brendanpshea.github.io}\\
    \faGithub\enspace\texttt{github.com/brendanpshea}\\
    \faEnvelope\enspace brendan.shea@rctc.edu%
  }}
\end{minipage}%
\hfill
\begin{minipage}[c]{0.155\linewidth}
  \centering
  \qrcode[height=2.45cm]{https://brendanpshea.github.io/computing_ai_ethics/}\\[3pt]
  {\sffamily\fontsize{5.0}{6.2}\selectfont\textcolor{white!80!mPrimary}%
    {brendanpshea.github.io/\linebreak computing\_ai\_ethics/}}
\end{minipage}
\end{tcolorbox}

% ---------------------------------------------------------------------------
%  PROJECT OVERVIEW  |  DISCIPLINE KEY  (two-panel row; QR is in header)
% ---------------------------------------------------------------------------
\begin{minipage}[t]{0.415\linewidth}
\begin{tcolorbox}[
  enhanced,
  colback=mAccent!8,
  colframe=mAccent,
  fonttitle=\small\sffamily\bfseries,
  coltitle=white,
  colbacktitle=mAccent!85!black,
  title={About This OER},
  boxrule=0.5pt,
  arc=2pt,
  left=6pt, right=6pt,
  top=4pt, bottom=4pt,
  toptitle=2pt, bottomtitle=2pt,
  before skip=0pt,
  after skip=4pt,
]
{\sffamily\footnotesize\textbf{Computing and AI Ethics} is a free,
openly-licensed textbook for any major. Twelve self-contained chapters
span philosophy through AI existential risk, each with editable lecture
slides, a full HTML chapter, and an auto-graded quiz. Every component
is on GitHub under \textbf{CC BY-NC-SA 4.0}. Adopt one chapter or
the whole course --- no permission needed.}
\end{tcolorbox}

\begin{tcolorbox}[
  enhanced,
  colback=mPrimary!5,
  colframe=mPrimary,
  fonttitle=\small\sffamily\bfseries,
  coltitle=white,
  colbacktitle=mPrimary,
  title={Discipline Tags},
  boxrule=0.5pt,
  arc=2pt,
  left=6pt, right=6pt,
  top=4pt, bottom=4pt,
  toptitle=2pt, bottomtitle=2pt,
  before skip=0pt,
  after skip=0pt,
]
{\sffamily\footnotesize
  \dtag{tagLA}{LA}\enspace Liberal Arts, Humanities, Social Sciences\\[2.5pt]
  \dtag{tagSCI}{SCI}\enspace Natural Sciences, Math, Environment\\[2.5pt]
  \dtag{tagHLTH}{HLTH}\enspace Health Sciences, Nursing, Allied Health\\[2.5pt]
  \dtag{tagBUS}{BUS}\enspace Business, Accounting, Economics\\[2.5pt]
  \dtag{tagIT}{IT}\enspace Information Technology, Computer Science\\[2.5pt]
  \dtag{tagEDU}{EDU}\enspace Education, Workforce Development
}
\end{tcolorbox}
\end{minipage}%
\hfill
\begin{minipage}[t]{0.565\linewidth}
\begin{tcolorbox}[
  enhanced,
  colback=mPrimary!5,
  colframe=mPrimary,
  fonttitle=\small\sffamily\bfseries,
  coltitle=white,
  colbacktitle=mPrimary,
  title={What Is Available for Each Chapter},
  boxrule=0.5pt,
  arc=2pt,
  left=6pt, right=6pt,
  top=4pt, bottom=4pt,
  toptitle=2pt, bottomtitle=2pt,
  before skip=0pt,
  after skip=4pt,
]
{\sffamily\footnotesize
\begin{tabular}{@{}lp{0.72\linewidth}@{}}
  \textbf{Lecture Slides} &
    Full Beamer/Metropolis deck in \LaTeX{} --- editable and remixable \\[2pt]
  \textbf{HTML Reading} &
    Complete chapter in accessible HTML; no login, no cost \\[2pt]
  \textbf{Quiz App} &
    Plain JavaScript, auto-graded, 10--20 Qs; host on GitHub Pages \\[2pt]
  \textbf{Source Code} &
    All \LaTeX{} and HTML source on GitHub --- fork freely\\
\end{tabular}
}
\end{tcolorbox}

\begin{tcolorbox}[
  enhanced,
  colback=white,
  colframe=mAccent,
  fonttitle=\small\sffamily\bfseries,
  coltitle=white,
  colbacktitle=mAccent!85!black,
  title={How to Use These Questions},
  boxrule=0.5pt,
  arc=2pt,
  left=6pt, right=6pt,
  top=4pt, bottom=4pt,
  toptitle=2pt, bottomtitle=2pt,
  before skip=0pt,
  after skip=0pt,
]
{\sffamily\footnotesize
\begin{tabular}{@{}lp{0.80\linewidth}@{}}
  \faEdit      & Assign as a short reflection connecting course content
                 to a current event or professional scenario.\\[2pt]
  \faComments  & Open class with a 2-min silent write, then discuss.\\[2pt]
  \faUsers     & Partner with a philosophy or CS instructor to
                 co-teach 1--2 chapters relevant to your field.\\[2pt]
  \faBookOpen  & Use as supplementary reading in any Gen Ed course ---
                 one chapter stands alone as a guest lecture or module.\\
\end{tabular}
}
\end{tcolorbox}
\end{minipage}

\vspace{5pt}
\noindent{\sffamily\bfseries\normalsize\textcolor{mPrimary}{%
  \faComments\enspace Questions for Cross-Curricular Teaching and Discussion}}\\[-1pt]
\noindent\textcolor{mAccent}{\rule{\linewidth}{1.3pt}}
\vspace{3pt}

% ---------------------------------------------------------------------------
%  TWELVE CHAPTER CARDS — 2-column layout
% ---------------------------------------------------------------------------
\setlist[itemize]{leftmargin=1em, itemsep=1.5pt, topsep=1pt, parsep=0pt}

\begin{multicols}{2}

% -- Ch 1 -------------------------------------------------------------------
\begin{tcolorbox}[chap,
  title={Ch.~1 \textnormal{---} History of IT Ethics: Plato to Web~1.0}]
\sffamily\footnotesize
\textit{Five revolutions---writing, print, telegraph, broadcast, internet---each
promised liberation while enabling new forms of control.}
\begin{itemize}
  \item Which groups historically controlled access to each new information
    technology --- and who was excluded? What parallels exist with AI adoption
    today? \dtag{tagLA}{LA} \dtag{tagSCI}{SCI}
  \item Plato feared that writing produces the \emph{appearance} of knowledge
    without genuine understanding. Does the same critique apply to search
    engines or large language models? \dtag{tagEDU}{EDU} \dtag{tagLA}{LA}
    \dtag{tagIT}{IT}
\end{itemize}
\end{tcolorbox}

% -- Ch 2 -------------------------------------------------------------------
\begin{tcolorbox}[chap,
  title={Ch.~2 \textnormal{---} Virtue Ethics and IT}]
\sffamily\footnotesize
\textit{Aristotle: virtues are habits formed by practice. Technology shapes what
we repeatedly do --- and thus who we become.}
\begin{itemize}
  \item Which habits of mind --- curiosity, empathy, patience, honesty ---
    does regular technology use in your field risk strengthening or eroding?
    \dtag{tagHLTH}{HLTH} \dtag{tagEDU}{EDU} \dtag{tagLA}{LA}
  \item If a practitioner delegates diagnosis, writing, or analysis to AI,
    are they still building the expertise their role requires --- or
    outsourcing it? \dtag{tagHLTH}{HLTH} \dtag{tagBUS}{BUS} \dtag{tagIT}{IT}
\end{itemize}
\end{tcolorbox}

% -- Ch 3 -------------------------------------------------------------------
\begin{tcolorbox}[chap,
  title={Ch.~3 \textnormal{---} Free Speech and Social Media}]
\sffamily\footnotesize
\textit{Platforms occupy new legal and ethical ground: neither public squares
nor private clubs, yet algorithmically shaping who is heard and why.}
\begin{itemize}
  \item If outrage drives engagement, who bears moral responsibility for
    political polarization --- users, platform designers, or the algorithms?
    \dtag{tagLA}{LA} \dtag{tagBUS}{BUS} \dtag{tagIT}{IT}
  \item Should healthcare workers, educators, or attorneys face special
    obligations about discussing clients or students on social media?
    \dtag{tagHLTH}{HLTH} \dtag{tagEDU}{EDU} \dtag{tagBUS}{BUS}
\end{itemize}
\end{tcolorbox}

% -- Ch 4 -------------------------------------------------------------------
\begin{tcolorbox}[chap,
  title={Ch.~4 \textnormal{---} Intellectual Property}]
\sffamily\footnotesize
\textit{AI trained on creative work can now produce outputs in the style of
living artists and writers --- without consent or compensation.}
\begin{itemize}
  \item Does Locke's principle --- that creators deserve the fruits of their
    labor --- support or undermine AI companies that train on existing work?
    \dtag{tagBUS}{BUS} \dtag{tagLA}{LA} \dtag{tagIT}{IT}
  \item When AI contributes to a scientific discovery or clinical innovation,
    who holds the patent --- and on what ethical or legal basis?
    \dtag{tagSCI}{SCI} \dtag{tagHLTH}{HLTH} \dtag{tagBUS}{BUS}
\end{itemize}
\end{tcolorbox}

% -- Ch 5 -------------------------------------------------------------------
\begin{tcolorbox}[chap,
  title={Ch.~5 \textnormal{---} Cryptocurrency}]
\sffamily\footnotesize
\textit{Cryptocurrency challenges state control of money and forces a basic
question: what makes a currency legitimate?}
\begin{itemize}
  \item Can a technology that simultaneously serves the unbanked poor and
    enables unprecedented-scale fraud be evaluated in simple ethical terms?
    \dtag{tagBUS}{BUS} \dtag{tagLA}{LA} \dtag{tagSCI}{SCI}
  \item Bitcoin uses as much energy annually as some medium-sized nations.
    How should environmental and social costs factor into financial and
    technology decisions? \dtag{tagSCI}{SCI} \dtag{tagBUS}{BUS}
    \dtag{tagEDU}{EDU}
\end{itemize}
\end{tcolorbox}

% -- Ch 6 -------------------------------------------------------------------
\begin{tcolorbox}[chap,
  title={Ch.~6 \textnormal{---} Privacy and Surveillance}]
\sffamily\footnotesize
\textit{Surveillance capitalism treats behavioral data as a raw material
to be harvested, predicted, and sold without users' meaningful awareness.}
\begin{itemize}
  \item ``If you have nothing to hide, you have nothing to fear.'' What are
    this claim's strongest weaknesses --- and which communities suffer most
    when it wins?
    \dtag{tagLA}{LA} \dtag{tagHLTH}{HLTH} \dtag{tagSCI}{SCI}
  \item Health records, fitness trackers, and insurance algorithms create rich
    behavioral profiles. Where should patient or consumer consent limits lie?
    \dtag{tagHLTH}{HLTH} \dtag{tagIT}{IT} \dtag{tagBUS}{BUS}
\end{itemize}
\end{tcolorbox}

% -- Ch 7 -------------------------------------------------------------------
\begin{tcolorbox}[chap,
  title={Ch.~7 \textnormal{---} Introduction to AI}]
\sffamily\footnotesize
\textit{AI has moved from explicit rules to neural networks trained on data ---
systems that sometimes outperform experts without being able to explain why.}
\begin{itemize}
  \item Searle's Chinese Room argues that symbol manipulation produces no real
    understanding. Does this matter for how much we should trust AI advice in
    high-stakes contexts? \dtag{tagLA}{LA} \dtag{tagSCI}{SCI} \dtag{tagIT}{IT}
  \item AI now exceeds human performance at several diagnostic tasks. What
    standards should govern AI adoption in clinical, legal, or financial
    settings --- and who is liable when it fails?
    \dtag{tagHLTH}{HLTH} \dtag{tagBUS}{BUS} \dtag{tagIT}{IT}
\end{itemize}
\end{tcolorbox}

% -- Ch 8 -------------------------------------------------------------------
\begin{tcolorbox}[chap,
  title={Ch.~8 \textnormal{---} AI and the Future of Work}]
\sffamily\footnotesize
\textit{Hannah Arendt distinguished \emph{labor} (biological necessity),
\emph{work} (skilled craft), and \emph{action} (civic life) --- AI encroaches
on all three, but unevenly.}
\begin{itemize}
  \item Which tasks in your field are most exposed to AI automation --- and
    which seem to require irreducibly human judgment, relationship, or
    creativity?
    \dtag{tagBUS}{BUS} \dtag{tagHLTH}{HLTH} \dtag{tagEDU}{EDU}
  \item If AI handles most routine cognitive work, what should higher education
    in your discipline \emph{become} --- and which skills should students be
    building right now? \dtag{tagEDU}{EDU} \dtag{tagBUS}{BUS} \dtag{tagIT}{IT}
\end{itemize}
\end{tcolorbox}

% -- Ch 9 -------------------------------------------------------------------
\begin{tcolorbox}[chap,
  title={Ch.~9 \textnormal{---} Social \& Environmental Impact of AI}]
\sffamily\footnotesize
\textit{AI's costs --- energy, water, e-waste, algorithmic bias --- are largely
invisible to users who benefit from its outputs.}
\begin{itemize}
  \item Training one large AI model can emit as much CO\textsubscript{2} as
    five cars over their lifetimes. How should we weigh environmental costs
    against AI's research, clinical, or business benefits?
    \dtag{tagSCI}{SCI} \dtag{tagBUS}{BUS} \dtag{tagHLTH}{HLTH}
  \item AI trained on historical data tends to amplify existing biases in
    hiring, lending, and medical diagnosis. What would genuinely fair AI
    look like in your profession --- and who gets to decide?
    \dtag{tagHLTH}{HLTH} \dtag{tagBUS}{BUS} \dtag{tagLA}{LA}
\end{itemize}
\end{tcolorbox}

% -- Ch 10 ------------------------------------------------------------------
\begin{tcolorbox}[chap,
  title={Ch.~10 \textnormal{---} AI and Existential Risk}]
\sffamily\footnotesize
\textit{Some researchers argue misaligned AI poses civilizational risk; critics
say this framing displaces focus from near-term harms happening right now.}
\begin{itemize}
  \item Should catastrophic-but-uncertain AI risks be governed like nuclear
    weapons --- with international treaties --- or does this framing distract
    from present injustices?
    \dtag{tagLA}{LA} \dtag{tagSCI}{SCI} \dtag{tagBUS}{BUS}
  \item How should scientists, journalists, and policymakers communicate about
    risks that are deeply uncertain but potentially severe? What makes
    risk communication responsible?
    \dtag{tagSCI}{SCI} \dtag{tagLA}{LA} \dtag{tagHLTH}{HLTH}
\end{itemize}
\end{tcolorbox}

% -- Ch 11 ------------------------------------------------------------------
\begin{tcolorbox}[chap,
  title={Ch.~11 \textnormal{---} Robot Rights}]
\sffamily\footnotesize
\textit{We infer consciousness in others by analogy to ourselves. AI challenges
this inference in ways philosophy and science have not yet resolved.}
\begin{itemize}
  \item When patients form genuine therapeutic bonds with care robots, do those
    systems acquire any moral weight --- or only the patients' interests
    matter? \dtag{tagHLTH}{HLTH} \dtag{tagLA}{LA} \dtag{tagSCI}{SCI}
  \item By what criteria --- sentience, rationality, relational bonds --- do we
    extend moral consideration to any being? Do current AI systems meet any of
    these? \dtag{tagSCI}{SCI} \dtag{tagLA}{LA} \dtag{tagIT}{IT}
\end{itemize}
\end{tcolorbox}

% -- Ch 12 ------------------------------------------------------------------
\begin{tcolorbox}[chap,
  title={Ch.~12 \textnormal{---} Games, Play, and Moral Philosophy}]
\sffamily\footnotesize
\textit{Bernard Suits argues that voluntary, rule-governed challenge-seeking is
the ideal human life --- a claim with radical implications if AI replaces most
necessary work.}
\begin{itemize}
  \item If AI handles most productive labor, is a life of games, art, and
    leisure genuinely fulfilling --- or does it hollow out human purpose and
    dignity? \dtag{tagLA}{LA} \dtag{tagEDU}{EDU} \dtag{tagBUS}{BUS}
  \item Loot boxes and engagement-maximizing design exploit psychological
    tendencies for profit. Are they ethically comparable to gambling --- and
    should they be regulated as such?
    \dtag{tagBUS}{BUS} \dtag{tagSCI}{SCI} \dtag{tagLA}{LA}
\end{itemize}
\end{tcolorbox}

\end{multicols}

% ---------------------------------------------------------------------------
%  CHAPTER ARC  (single-line summary at bottom)
% ---------------------------------------------------------------------------
\vspace{1pt}
\noindent\textcolor{mAccent}{\rule{\linewidth}{0.8pt}}
\vspace{2pt}

\begin{tcolorbox}[
  enhanced,
  colback=mPrimary!5,
  colframe=mPrimary,
  boxrule=0.5pt,
  arc=2pt,
  left=7pt, right=7pt,
  top=4pt, bottom=4pt,
  before skip=0pt,
  after skip=0pt,
]
{\sffamily\footnotesize\textcolor{mPrimary}{%
\textbf{Course arc:}
Chapters 1--2 build \textbf{philosophical foundations} (virtue ethics, the history of IT ethics).
Chapters 3--6 engage \textbf{pressing policy debates} (speech, IP, crypto, privacy).
Chapters 7--12 focus on \textbf{AI specifically} --- from what AI is through existential risk,
robot rights, and the ethics of play. Any chapter stands alone as a one-week module,
discussion unit, or guest lecture. All materials at
\href{https://brendanpshea.github.io/computing_ai_ethics/}%
     {\textbf{brendanpshea.github.io/computing\_ai\_ethics/}} ---
fork, adapt, and use freely.
}}
\end{tcolorbox}

\end{document}
